\chapter{Suprapubic Catheterization}

\section*{What Is This Test or Procedure?}
Suprapubic catheterization is the placement of a drainage catheter into the bladder through a puncture in the lower abdominal wall, above the pubic symphysis. It provides bladder drainage when urethral catheterization is impossible, contraindicated, or undesirable.
\section*{Alternative Names}
Suprapubic cystostomy; suprapubic bladder drainage; punch suprapubic catheter insertion
\section*{Principle}
With the bladder distended, a catheter is passed through the skin, rectus sheath, and bladder wall into the bladder lumen. The inflated balloon holds the catheter in place, and drainage is achieved by gravity or intermittent clamping.
\section*{History}
Suprapubic bladder drainage was practiced in antiquity as a treatment for urinary retention. The modern method, using a percutaneous trocar and a self-retaining catheter, was popularized in the mid-20th century as a safer alternative to repeated urethral instrumentation.
\section*{Why This Test or Procedure Is Ordered}
Ordered for urinary retention when urethral catheterization fails or is contraindicated, for long-term bladder drainage in men to reduce urethral complications, and for patients with urethral injury, prostatic surgery, or recurrent urethral catheter problems.
\section*{Is This a Primary or Secondary Test or Procedure}
Primary therapeutic procedure for bladder drainage when the urethral route is unsuitable.
\section*{How This Test or Procedure Is Performed}
The lower abdomen is prepared, and the bladder is confirmed distended by palpation or ultrasound. The puncture site is marked 2-3 cm above the pubic symphysis and anesthetized. A trocar and catheter are advanced through the abdominal wall into the bladder, the catheter is advanced, the trocar is removed, and the balloon is inflated and secured to the abdominal wall.
\section*{What Preparation Is Needed}
Informed consent, confirmation of a distended bladder, and imaging guidance when the bladder cannot be clearly palpated or there is doubt about anatomy. Anticoagulant use and bleeding risk are reviewed.
\section*{Contraindications and Precautions}
Contraindications include an empty or non-distended bladder, severe bleeding diathesis, known bladder cancer or previous lower abdominal surgery (potential bowel adhesions), and skin infection at the site. Caution is required if bowel overlies the puncture path.
\section*{Risks}
Hematuria, infection, bowel injury or perforation, puncture of abdominal vessels, catheter displacement, and leakage around the catheter site.
\section*{What Happens After the Test or Procedure}
The catheter is connected to a drainage bag and the site is kept clean and dry. The urine is observed for hematuria and the patient is monitored for abdominal pain or signs of infection.
\section*{Understanding Your Results}
A free flow of urine confirms correct placement, and clear urine with no abdominal tenderness suggests an uncomplicated procedure. Difficult insertion, blood, or fever may indicate complications.
\section*{Normal Results}
A patent suprapubic catheter draining clear urine with a clean, non-inflamed exit site and no leakage or infection.
\section*{Follow-up Tests and Procedures}
The catheter site and output are monitored, the catheter is changed at scheduled intervals, and a urinalysis or culture is performed if infection is suspected. Trial of clamping or voiding is attempted before eventual removal.
\section*{References}
\begin{enumerate}
  \item MedlinePlus. Suprapubic catheter care. U.S. National Library of Medicine; 2025.
  \item Current Medical Diagnosis and Treatment. McGraw-Hill; 2025.
\end{enumerate}

\clearpage
